\documentclass[12pt,letterpaper]{article}
\usepackage{graphicx,textcomp}
\usepackage{natbib}
\usepackage{setspace}
\usepackage{fullpage}
\usepackage{color}
\usepackage[reqno]{amsmath}
\usepackage{amsthm}
\usepackage{fancyvrb}
\usepackage{amssymb,enumerate}
\usepackage[all]{xy}
\usepackage{endnotes}
\usepackage{lscape}
\newtheorem{com}{Comment}
\usepackage{float}
\usepackage{hyperref}
\newtheorem{lem} {Lemma}
\newtheorem{prop}{Proposition}
\newtheorem{thm}{Theorem}
\newtheorem{defn}{Definition}
\newtheorem{cor}{Corollary}
\newtheorem{obs}{Observation}
\usepackage[compact]{titlesec}
\usepackage{dcolumn}
\usepackage{tikz}
\usetikzlibrary{arrows}
\usepackage{multirow}
\usepackage{xcolor}
\newcolumntype{.}{D{.}{.}{-1}}
\newcolumntype{d}[1]{D{.}{.}{#1}}
\definecolor{light-gray}{gray}{0.65}
\usepackage{url}
\usepackage{listings}
\usepackage{color}
\usepackage{relsize}
\usepackage{graphicx}
\usepackage{amssymb}


\definecolor{codegreen}{rgb}{0,0.6,0}
\definecolor{codegray}{rgb}{0.5,0.5,0.5}
\definecolor{codepurple}{rgb}{0.58,0,0.82}
\definecolor{backcolour}{rgb}{0.95,0.95,0.92}

\lstdefinestyle{mystyle}{
	backgroundcolor=\color{backcolour},   
	commentstyle=\color{codegreen},
	keywordstyle=\color{magenta},
	numberstyle=\tiny\color{codegray},
	stringstyle=\color{codepurple},
	basicstyle=\footnotesize,
	breakatwhitespace=false,         
	breaklines=true,                 
	captionpos=b,                    
	keepspaces=true,                 
	numbers=left,                    
	numbersep=5pt,                  
	showspaces=false,                
	showstringspaces=false,
	showtabs=false,                  
	tabsize=2
}
\lstset{style=mystyle}
\newcommand{\Sref}[1]{Section~\ref{#1}}
\newtheorem{hyp}{Hypothesis}

\title{Problem Set 3}
\date{Due: November 13, 2025}
\author{Rosalie LaCroix}


\begin{document}
	\maketitle
	\section*{Instructions}
	\begin{itemize}
		\item Please show your work! You may lose points by simply writing in the answer. If the problem requires you to execute commands in \texttt{R}, please include the code you used to get your answers. Please also include the \texttt{.R} file that contains your code. If you are not sure if work needs to be shown for a particular problem, please ask.
		\item Your homework should be submitted electronically on GitHub.
		\item This problem set is due before 23:59 on Thursday November 13, 2025. No late assignments will be accepted.
		
	\end{itemize}
	
	\vspace{.25cm}
	
	\noindent In this problem set, you will run several regressions and create an add variable plot (see the lecture slides) in \texttt{R} using the \texttt{incumbents\_subset.csv} dataset. Include all of your code.
	
	\vspace{.5cm}
	\section*{Question 1}
	\vspace{.25cm}
	\noindent We are interested in knowing how the difference in campaign spending between incumbent and challenger affects the incumbent's vote share. 
	\begin{enumerate}
		\item Run a regression where the outcome variable is \texttt{voteshare} and the explanatory variable is \texttt{difflog}.
		\lstinputlisting[language=R, firstline=40, lastline=40]{PS03_myanswers_RL.R}
		\begin{table}[!htbp] \centering 
			\caption{Table of Coefficients} 
			\label{} 
			\begin{tabular}{@{\extracolsep{5pt}}lc} 
				\\[-1.8ex]\hline 
				\hline \\[-1.8ex] 
				& \multicolumn{1}{c}{\textit{Dependent variable:}} \\ 
				\cline{2-2} 
				\\[-1.8ex] & voteshare \\ 
				& Regression 1 \\ 
				\hline \\[-1.8ex] 
				difflog & 0.042$^{***}$ \\ 
				& (0.001) \\ 
				& \\ 
				Constant & 0.579$^{***}$ \\ 
				& (0.002) \\ 
				& \\ 
				\hline \\[-1.8ex] 
				Observations & 3,193 \\ 
				R$^{2}$ & 0.367 \\ 
				Adjusted R$^{2}$ & 0.367 \\ 
				Residual Std. Error & 0.079 (df = 3191) \\ 
				F Statistic & 1,852.791$^{***}$ (df = 1; 3191) \\ 
				\hline 
				\hline \\[-1.8ex] 
				\textit{Note:}  & \multicolumn{1}{r}{$^{*}$p$<$0.1; $^{**}$p$<$0.05; $^{***}$p$<$0.01} \\ 
			\end{tabular} 
		\end{table} 
		
		\newpage
		\item Make a scatterplot of the two variables and add the regression line. 
		\lstinputlisting[language=R, firstline=45, lastline=49]{PS03_myanswers_RL.R}
		\includegraphics[width=10cm]{scatterplot_1.pdf}
		\item Save the residuals of the model in a separate object.	
		\lstinputlisting[language=R, firstline=53, lastline=53]{PS03_myanswers_RL.R}
		\item Write the prediction equation.
		\lstinputlisting[language=R, firstline=56, lastline=56]{PS03_myanswers_RL.R}
		
		\textsmaller[1]{The prediction equation is y = 0.5790 + 0.0417x or voteshare = 0.5790 + 0.0417difflog in the context of this specific data.}
		
	\end{enumerate}
	
	\newpage
	
	\section*{Question 2}
	\noindent We are interested in knowing how the difference between incumbent and challenger's spending and the vote share of the presidential candidate of the incumbent's party are related.	\vspace{.25cm}
	\begin{enumerate}
		\item Run a regression where the outcome variable is \texttt{presvote} and the explanatory variable is \texttt{difflog}.	
		\lstinputlisting[language=R, firstline=62, lastline=62]{PS03_myanswers_RL.R}
		\begin{table}[!htbp] \centering 
			\caption{Table of Coefficients} 
			\label{} 
			\begin{tabular}{@{\extracolsep{5pt}}lc} 
				\\[-1.8ex]\hline 
				\hline \\[-1.8ex] 
				& \multicolumn{1}{c}{\textit{Dependent variable:}} \\ 
				\cline{2-2} 
				\\[-1.8ex] & presvote \\ 
				& Regression 2 \\ 
				\hline \\[-1.8ex] 
				difflog & 0.024$^{***}$ \\ 
				& (0.001) \\ 
				& \\ 
				Constant & 0.508$^{***}$ \\ 
				& (0.003) \\ 
				& \\ 
				\hline \\[-1.8ex] 
				Observations & 3,193 \\ 
				R$^{2}$ & 0.088 \\ 
				Adjusted R$^{2}$ & 0.088 \\ 
				Residual Std. Error & 0.110 (df = 3191) \\ 
				F Statistic & 307.715$^{***}$ (df = 1; 3191) \\ 
				\hline 
				\hline \\[-1.8ex] 
				\textit{Note:}  & \multicolumn{1}{r}{$^{*}$p$<$0.1; $^{**}$p$<$0.05; $^{***}$p$<$0.01} \\ 
			\end{tabular} 
		\end{table}
		\item Make a scatterplot of the two variables and add the regression line. 
		\lstinputlisting[language=R, firstline=67, lastline=71]{PS03_myanswers_RL.R}
		\includegraphics[width=10cm]{scatterplot_2.pdf}
		\item Save the residuals of the model in a separate object.	
		\lstinputlisting[language=R, firstline=75, lastline=75]{PS03_myanswers_RL.R}
		
		\newpage
		\item Write the prediction equation.
		\lstinputlisting[language=R, firstline=78, lastline=78]{PS03_myanswers_RL.R}
		
		\textsmaller[1]{The prediction equation is y = 0.5076 + 0.0238x or presvote = 0.5076 + 0.0238difflog in the context of this specific data.}
		
	\end{enumerate}
	
	\newpage	
	\section*{Question 3}
	
	\noindent We are interested in knowing how the vote share of the presidential candidate of the incumbent's party is associated with the incumbent's electoral success.
	\vspace{.25cm}
	\begin{enumerate}
		\item Run a regression where the outcome variable is \texttt{voteshare} and the explanatory variable is \texttt{presvote}.
		\lstinputlisting[language=R, firstline=84, lastline=84]{PS03_myanswers_RL.R}
		\begin{table}[!htbp] \centering 
			\caption{Table of Coefficients} 
			\label{} 
			\begin{tabular}{@{\extracolsep{5pt}}lc} 
				\\[-1.8ex]\hline 
				\hline \\[-1.8ex] 
				& \multicolumn{1}{c}{\textit{Dependent variable:}} \\ 
				\cline{2-2} 
				\\[-1.8ex] & voteshare \\ 
				& Regression 3 \\ 
				\hline \\[-1.8ex] 
				presvote & 0.388$^{***}$ \\ 
				& (0.013) \\ 
				& \\ 
				Constant & 0.441$^{***}$ \\ 
				& (0.008) \\ 
				& \\ 
				\hline \\[-1.8ex] 
				Observations & 3,193 \\ 
				R$^{2}$ & 0.206 \\ 
				Adjusted R$^{2}$ & 0.206 \\ 
				Residual Std. Error & 0.088 (df = 3191) \\ 
				F Statistic & 826.950$^{***}$ (df = 1; 3191) \\ 
				\hline 
				\hline \\[-1.8ex] 
				\textit{Note:}  & \multicolumn{1}{r}{$^{*}$p$<$0.1; $^{**}$p$<$0.05; $^{***}$p$<$0.01} \\ 
			\end{tabular} 
		\end{table}
		\item Make a scatterplot of the two variables and add the regression line. 
		\lstinputlisting[language=R, firstline=89, lastline=93]{PS03_myanswers_RL.R}
		\includegraphics[width=10cm]{scatterplot_3.pdf}
		
		\newpage
		\item Write the prediction equation.
		\lstinputlisting[language=R, firstline=97, lastline=97]{PS03_myanswers_RL.R}
		
		\textsmaller[1]{The prediction equation is y = 0.4413 + 0.3880x or voteshare = 0.4413 + 0.3880presvote in the context of this specific data.}
		
	\end{enumerate}
	
	
	\newpage	
	\section*{Question 4}
	\noindent The residuals from part (a) tell us how much of the variation in \texttt{voteshare} is $not$ explained by the difference in spending between incumbent and challenger. The residuals in part (b) tell us how much of the variation in \texttt{presvote} is $not$ explained by the difference in spending between incumbent and challenger in the district.
	\begin{enumerate}
		\item Run a regression where the outcome variable is the residuals from Question 1 and the explanatory variable is the residuals from Question 2.
		\lstinputlisting[language=R, firstline=103, lastline=103]{PS03_myanswers_RL.R}
		\begin{table}[!htbp] \centering 
			\caption{Table of Coefficients} 
			\label{} 
			\begin{tabular}{@{\extracolsep{5pt}}lc} 
				\\[-1.8ex]\hline 
				\hline \\[-1.8ex] 
				& \multicolumn{1}{c}{\textit{Dependent variable:}} \\ 
				\cline{2-2} 
				\\[-1.8ex] & residuals1 \\ 
				& Residual Regression \\ 
				\hline \\[-1.8ex] 
				residuals2 & 0.257$^{***}$ \\ 
				& (0.012) \\ 
				& \\ 
				Constant & $-$0.000 \\ 
				& (0.001) \\ 
				& \\ 
				\hline \\[-1.8ex] 
				Observations & 3,193 \\ 
				R$^{2}$ & 0.130 \\ 
				Adjusted R$^{2}$ & 0.130 \\ 
				Residual Std. Error & 0.073 (df = 3191) \\ 
				F Statistic & 476.975$^{***}$ (df = 1; 3191) \\ 
				\hline 
				\hline \\[-1.8ex] 
				\textit{Note:}  & \multicolumn{1}{r}{$^{*}$p$<$0.1; $^{**}$p$<$0.05; $^{***}$p$<$0.01} \\ 
			\end{tabular} 
		\end{table}
		\item Make a scatterplot of the two residuals and add the regression line. 
		\lstinputlisting[language=R, firstline=107, lastline=113]{PS03_myanswers_RL.R}
		\includegraphics[width=10cm]{scatterplot_4.pdf}
		\item Write the prediction equation.
		
		\textsmaller[1]{The prediction equation is y = -1.942e-18 + 0.2569x or residuals1 = -1.942e-18 + 0.2569residuals2 in the context of this specific data.}
	\end{enumerate}
	
	\newpage	
	
	\section*{Question 5}
	\noindent What if the incumbent's vote share is affected by both the president's popularity and the difference in spending between incumbent and challenger? 
	\begin{enumerate}
		\item Run a regression where the outcome variable is the incumbent's \texttt{voteshare} and the explanatory variables are \texttt{difflog} and \texttt{presvote}.
		\lstinputlisting[language=R, firstline= 123, lastline= 123]{PS03_myanswers_RL.R}
		\begin{table}[!htbp] \centering 
			\caption{Table of Coefficients} 
			\label{} 
			\begin{tabular}{@{\extracolsep{5pt}}lc} 
				\\[-1.8ex]\hline 
				\hline \\[-1.8ex] 
				& \multicolumn{1}{c}{\textit{Dependent variable:}} \\ 
				\cline{2-2} 
				\\[-1.8ex] & voteshare \\ 
				& Regression 4 \\ 
				\hline \\[-1.8ex] 
				difflog & 0.036$^{***}$ \\ 
				& (0.001) \\ 
				& \\ 
				presvote & 0.257$^{***}$ \\ 
				& (0.012) \\ 
				& \\ 
				Constant & 0.449$^{***}$ \\ 
				& (0.006) \\ 
				& \\ 
				\hline \\[-1.8ex] 
				Observations & 3,193 \\ 
				R$^{2}$ & 0.450 \\ 
				Adjusted R$^{2}$ & 0.449 \\ 
				Residual Std. Error & 0.073 (df = 3190) \\ 
				F Statistic & 1,302.947$^{***}$ (df = 2; 3190) \\ 
				\hline 
				\hline \\[-1.8ex] 
				\textit{Note:}  & \multicolumn{1}{r}{$^{*}$p$<$0.1; $^{**}$p$<$0.05; $^{***}$p$<$0.01} \\ 
			\end{tabular} 
		\end{table} 
		\item Write the prediction equation.
		\lstinputlisting[language=R, firstline=127, lastline=127]{PS03_myanswers_RL.R}
		\textsmaller[1]{The prediction equation is y = 0.4486 + 0.0355$x_{1}$ + 0.2569$x_{2}$ or voteshare = 0.4486 + 0.0355difflog + 0.2569presvote in the context of this specific regression.}
		\item What is it in this output that is identical to the output in Question 4? Why do you think this is the case?
		
		\textsmaller[1]{In this output and in the output in Question 4, the slope for presvote impact on voteshare is identical to the slope for residuals2 impact on of residuals1, which are both 0.2569. These two values are the same as they both represent the effect of the vote share of the presidental candidate of the incumbent's party on the incumbent's vote share whilst controlling for the confounding effect of the difference in campaign spending between incumbent and challenger. This is due to the fact that the residuals from the regressions in question 1 and 2 (residuals1 and 2, respectively), represent the incumbent's vote share and vote share of the presidential candidate of the incumbent's party that is not explained by the difference in campaign spending between incumbent and challenger. When these two residuals are ran in a regression with one another, the slope of residuals2 effect on residuals1 represents the effect of presvote on voteshare controlling for difflog. Running a multiple regression of the effects of both presvote and difflog on voteshare, the coefficient for presvote will therefore be the same as residuals2 coefficient as multiple regression coefficients statistically control for the effects of other variables.}
		
		
	\end{enumerate}
	
	
	
	
\end{document}
